\chapter{離婁章句下}
\kaishi*

\section{論先聖後聖}

孟子曰：「舜生於諸馮，遷於負夏，卒於鳴條，東夷之人也。」\jiazhu[1]{生始，卒終，記終始也。諸馮、負夏、鳴條，皆地名。負，負海也，在東方夷服之地，故曰東夷之人也。}

「文王生於岐周，卒於畢郢，西夷之人也。」\jiazhu[1]{岐周、畢郢，地名也。岐山下周之舊邑，近畎夷。畎夷在西，故曰西夷之人也。《書》曰：「太子發上祭于畢，下至于盟津。」畢，文王墓，近於酆鎬也。}

「地之相去也，千有餘里；世之相後也，千有餘歲。得志行乎中國，若合符節。先聖後聖，其揆一也。」\jiazhu[1]{土地相去千有餘里以外也。舜至文王千二百歲。得志行政於中國，謂王也。如合符節。節，玉節也。《周禮》有六節。揆，度也。言聖人之度量同也。}\par \jiazhu[1]{章指言：聖人殊世而合其道，地雖不比，由通一軌，故可以爲百王法也。}

\section{論子産惠而不知爲政}

子産聽鄭國之政，以其乘輿濟人於溱洧。\jiazhu[1]{子産，鄭卿，爲政聽訟也。溱、洧，水名。見人有冬涉者，仁心不忍，以其乘車渡之也。}

孟子曰：「惠而不知爲政。歲十一月，徒杠成；十二月，輿梁成，民未病涉也。」\jiazhu[1]{以爲子産有惠民之用，而不知爲政當以時修橋梁。民何由病苦涉水乎？周十一月，夏九月，可以成步渡之功；周十二月，夏十月，可以成輿梁也。}

「君子平其政，行辟人可也，焉得人人而濟之？故爲政者，每人而悅之，日亦不足矣。」\jiazhu[1]{君子爲國家平治政事刑法，使無違失，其道辟除人，使卑辟尊，可爲也，安得人人濟渡於水乎？每人輒欲自加恩以悅其意，則日力不足以足之也。}\par \jiazhu[1]{章指言：重民之道，平政爲首。人君由天，天不家撫，是故子産渡人，孟子不取也。}

\section{論君臣相待之道}

孟子告齊宣王曰：「君之視臣如手足，則臣視君如腹心；君之視臣如犬馬，則臣視君如國人；君之視臣如土芥，則臣視君如寇讎。」\jiazhu[1]{芥，草芥也。臣緣君恩以爲差等，其心所執若是也。}

王曰：「禮爲舊君有服，何如斯可爲服矣？」\jiazhu[1]{宣王問禮舊臣爲舊君服喪，問君何如則可以爲之服也。}

曰：「諫行言聽，膏澤下於民；有故而去，則君使人導之出疆，又先於其所往；去三年不反，然後收其田里。此之謂三有禮焉。如此，則爲之服矣。」\jiazhu[1]{諫行，從也。膏澤，恩惠也。有故，以他事去也。導之出疆，防剽掠也。先於其所往，稱其德也。三年不反，乃收其田禄及里居也。此三者有禮，則爲之服矣。}

「今也爲臣，諫則不行，言則不聽，膏澤不下於民；有故而去，則君搏執之，又極之於其所往；去之日，遂收其田里。此之謂寇讎。寇讎何服之有？」\jiazhu[1]{搏執其族親也。極者，惡而困之也。遇臣若寇讎，何服之有乎？}\par \jiazhu[1]{章指言：君臣之道，以義爲表，以恩爲裏，表裏相應，猶若影響。舊君之服，蓋有所興，諷諭宣王，勸以仁也。}

\section{論士大夫去就}

孟子曰：「無罪而殺士，則大夫可以去；無罪而戮民，則士可以徙。」\jiazhu[1]{惡傷其類，視其下等，懼次及也。語曰：「鳶鵲蒙害，仁鳥曾逝。」此之謂也。}\par \jiazhu[1]{章指言：君子見幾而作，故趙殺鳴犢，孔子臨河而不濟也。}

\section{論君仁莫不仁}

孟子曰：「君仁莫不仁，君義莫不義。」\jiazhu[1]{君者，一國所瞻仰以爲法，故必從之。}\par \jiazhu[1]{章指言：君以仁義率衆，孰不順焉？上爲下效也。}

\section{論非禮非義}

孟子曰：「非禮之禮，非義之義，大人弗爲。」\jiazhu[1]{若禮而非禮，陳質娶婦而長拜之也；若義而非義，藉交報仇是也。此皆大人所不爲也。}\par \jiazhu[1]{章指言：禮義，人之所以折中，履其正者，乃可爲中。是以大人不行疑禮。}

\section{論中養不中}

孟子曰：「中也養不中，才也養不才，故人樂有賢父兄也。」\jiazhu[1]{中者，履中和之氣所生，謂之賢；才者，謂人之有俊才者。有此賢者，當以養育教誨不能，進之以善，故樂父兄之賢以養己也。}

「如中也棄不中，才也棄不才，則賢不肖之相去，其間不能以寸。」\jiazhu[1]{如使賢者棄愚，不養其所當養，則賢亦近愚矣。如此，賢不肖相覺，何能分寸？明不可不相訓導也。}\par \jiazhu[1]{章指言：父兄已賢，子弟既頑，教而不改，乃歸自然。}

\section{論有爲有守}

孟子曰：「人有不爲也，而後可以有爲。」\jiazhu[1]{人不爲苟得，乃能有讓千乘之志。}\par \jiazhu[1]{章指言：貴廉賤恥，乃有不爲。不爲非義，義乃可申。}

\section{論言人之不善}

孟子曰：「言人之不善，當如後患何？」\jiazhu[1]{人之有惡，惡人言之，言之當如後有患難及己乎？}\par \jiazhu[1]{章指言：好言人惡，殆非君子。故曰：「不忮不求，何用不臧。」}

\section{論仲尼不爲已甚}

孟子曰：「仲尼不爲已甚者。」\jiazhu[1]{仲尼彈邪以正，正斯可矣，不欲其已甚泰過也。}\par \jiazhu[1]{章指言：論曰：「疾之已甚，亂也。」故孟子譏踰牆距門者也。}

\section{論大人之行}

孟子曰：「大人者，言不必信，行不必果，惟義所在。」\jiazhu[1]{果，能也。大人杖義，義有不得必信其言，子爲父隱也；有不能得果行其所欲行者，若親在不得以其身許友也。義或重於信，故曰惟義所在。}\par \jiazhu[1]{章指言：大人之行，行其重者，不信不果，所求合義也。}

\section{論赤子之心}

孟子曰：「大人者，不失其赤子之心者也。」\jiazhu[1]{大人謂君。國君視民當如赤子，不失其民心之謂也。一說曰：赤子，嬰兒也。少小之心，專一未變化。人能不失其赤子時心，則爲貞正大人也。}\par \jiazhu[1]{章指言：人之所愛，莫過赤子，視民則然，民懷之矣。大人之行，不過是也。}

\section{論養生送死}

孟子曰：「養生者不足以當大事，惟送死可以當大事。」\jiazhu[1]{孝子事親致養，未足以爲大事；送終如禮，則爲能奉大事也。}\par \jiazhu[1]{章指言：養生竭力，人情所勉；哀死送終，行之高者。事不違禮，可謂難矣，故謂之大事。}

\section{論深造自得}

孟子曰：「君子深造之以道，欲其自得之也。」\jiazhu[1]{造，致也。言君子學問之法，欲深致極竟之，以知道意，欲使己得其原本，如性自有之也。}

「自得之，則居之安；居之安，則資之深；資之深，則取之左右逢其原。故君子欲其自得之也。」\jiazhu[1]{居之安，若己所自有也。資，取也。取之深，則得其根也。左右取之，在所逢遇，皆知其原本也。故使君子欲自得之也。}\par \jiazhu[1]{章指言：學必根原，如性自得，物來能名，事來不惑。君子好之，朝益暮習，道所以臻也。}

\section{論博學反約}

孟子曰：「博學而詳說之，將以反說約也。」\jiazhu[1]{博，廣；詳，悉也。廣學悉其微言而說之者，將以約說其要。意不盡知，則不能要言之也。}\par \jiazhu[1]{章指言：廣尋道意，詳說其事，要約至義，還反於樸，說之美者也。}

\section{論以善養人}

孟子曰：「以善服人者，未有能服人者也；以善養人，然後能服天下。天下不心服而王者，未之有也。」\jiazhu[1]{以善服人之道治世，謂以威力服人者也，故人不心服。以善養人，養之以仁恩，然後心服矣。文王治岐是也。天下不心服，何由而王也？}\par \jiazhu[1]{章指言：五伯服人，三王服心，其服一也，功則不同。上論堯舜，其是違乎？}

\section{論言無實不祥}

孟子曰：「言無實不祥。不祥之實，蔽賢者當之。」\jiazhu[1]{凡言皆有實。孝子之實，養親是也；善之實，仁義是也。祥，善；當，直也。不善之實，何等也？蔽賢之人，直於不善之實也。}\par \jiazhu[1]{章指言：進賢受上賞，蔽賢蒙顯戮，故謂之不祥也。}

\section{徐子問仲尼稱水}

徐子曰：「仲尼亟稱於水，曰『水哉，水哉！』何取於水也？」\jiazhu[1]{徐子，徐辟也。問仲尼何取於水而稱之也？}

孟子曰：「原泉混混，不舍晝夜，盈科而後進，放乎四海。有本者如是，是之取爾。」\jiazhu[1]{言水不舍晝夜而進，盈滿科坎，放，至也。至於四海者，有原本也。以況於事有本者皆如是，是之取也。}

「苟爲無本，七八月之間雨集，溝澮皆盈，其涸也可立而待也。」\jiazhu[1]{苟，誠也。誠令無本，若周七八月，夏五六月，天之大雨，潦水卒集，大溝小澮皆滿，然其涸乾可立待者，無本之故也。}

「故聲聞過情，君子恥之。」\jiazhu[1]{人無本行，暴得善聲，令聞過其情，若潦水不能久也，故君子恥之。}\par \jiazhu[1]{章指言：有本不竭，無本則涸。虛聲過實，君子恥諸。是以仲尼在川上曰：「逝者如斯。」}

\section{論人禽之辨}

孟子曰：「人之所以異於禽獸者幾希，庶民去之，君子存之。」\jiazhu[1]{幾希，無幾也。知義與不知義之間耳。衆民去義，君子存義也。}

「舜明於庶物，察於人倫，由仁義行，非行仁義也。」\jiazhu[1]{倫，序；察，識也。舜明庶物之情，識人事之序。仁義生於內，由其中而行，非強力行仁義也。故道性善，言必稱堯舜。}\par \jiazhu[1]{章指言：人與禽獸俱含天氣，就利辟害，其間不希。衆人皆然，君子則否。聖人超絕，識仁義之生於己也。}

\section{論禹、湯、文、武、周公}

孟子曰：「禹惡旨酒而好善言。」\jiazhu[1]{旨酒，美酒也。儀狄作酒，禹飲而甘之，遂疏儀狄而絕旨酒。《書》曰：「禹拜讜言。」}

「湯執中，立賢無方。」\jiazhu[1]{執中正之道，惟賢速立之，不問其從何方來。舉伊尹以爲相也。}

「文王視民如傷，望道而未之見。」\jiazhu[1]{視民如傷者，雍容不動擾也。望道而未至，殷録未盡，尚有賢臣，道未得至，故望而不致誅於紂也。}

「武王不泄邇，不忘遠。」\jiazhu[1]{泄，狎；邇，近也。不泄狎近賢，不遺忘遠善。近謂朝臣，遠謂諸侯也。}

「周公思兼三王，以施四事。其有不合者，仰而思之，夜以繼日；幸而得之，坐以待旦。」\jiazhu[1]{三王，三代之王也。四事，禹、湯、文、武所行事也。不合，已行有不合世，仰而思之，參諸天也。坐而待旦，言欲急施之也。}\par \jiazhu[1]{章指言：周公能思三王之道，以輔成王太平之隆，禮樂之備，蓋由此也。}

\section{論詩亡春秋作}

孟子曰：「王者之迹熄而詩亡，詩亡然後《春秋》作。」\jiazhu[1]{王者，謂聖王也。太平道衰，王迹止熄，頌聲不作，故詩亡。《春秋》撥亂，作於衰世也。}

「晉之《乘》，楚之《檮杌》，魯之《春秋》，一也。其事則齊桓、晉文，其文則史。孔子曰：『其義則丘竊取之矣。』」\jiazhu[1]{此三大國史記之名異。《乘》者，興於田賦乘馬之事，因以爲名；《檮杌》者，嚚凶之類，興於記惡之戒，因以爲名；《春秋》，以二始舉四時，記萬事之名。其事則五伯所理也。桓、文，五伯之盛者，故舉之。其文，史記之文也。孔子自謂竊取之，以爲素王也。孔子人臣，不受君命，私作之，故言竊，亦聖人之謙辭。}\par \jiazhu[1]{章指言：詩可以言，頌詠太平，時無所詠，《春秋》乃興。假史記之文，孔子正之以匡邪也。}

\section{論君子小人之澤}

孟子曰：「君子之澤，五世而斬；小人之澤，五世而斬。」\jiazhu[1]{澤者，滋潤之澤。大德大凶，流及後世。自高祖至玄孫，善惡之氣乃斷，故曰五世而斬。}

「予未得爲孔子徒也，予私淑諸人也。」\jiazhu[1]{予，我也。我未得爲孔子門徒也。淑，善也。我私善之於賢人耳，恨不得學於大聖也。}\par \jiazhu[1]{章指言：五世一體，上下通流，君子小人，斬各有時。企以高山，跌以陷汙，是以孟子恨不及乎仲尼也。}

\section{論取與死之傷}

孟子曰：「可以取，可以無取，取傷廉；可以與，可以無與，與傷惠；可以死，可以無死，死傷勇。」\jiazhu[1]{三者皆謂事可出入，不至違義，但傷此名，亦不陷於惡也。}\par \jiazhu[1]{章指言：廉、惠、勇，人之高行也。喪此三名，烈士病諸。故設斯科，以進能者也。}

\section{論羿之罪}

逢蒙學射於羿，盡羿之道，思天下惟羿爲愈己，於是殺羿。\jiazhu[1]{羿，有窮后羿。逢蒙，羿之家衆也。《春秋傳》曰：「羿將歸自田，家衆殺之。」}

孟子曰：「是亦羿有罪焉。」\jiazhu[1]{罪羿不擇人也。故以下事喻之。}

公明儀曰：「宜若無罪焉？」曰：「薄乎云爾，惡得無罪？」\jiazhu[1]{公明儀以爲羿無罪。孟子以爲薄於罪耳，安得無罪乎？}

「鄭人使子濯孺子侵衛，衛使庾公之斯追之。子濯孺子曰：『今日我疾作，不可以執弓，吾死矣夫！』」\jiazhu[1]{孺子，鄭大夫；庾公，衛大夫。疾作，瘧疾。}

「問其僕曰：『追我者誰也？』其僕曰：『庾公之斯也。』曰：『吾生矣。』」\jiazhu[1]{僕，御也。孺子曰吾必生矣。}

「其僕曰：『庾公之斯，衛之善射者也。夫子曰「吾生」，何謂也？』曰：『庾公之斯學射於尹公之他，尹公之他學射於我。夫尹公之他，端人也，其取友必端矣。』」\jiazhu[1]{端人，用心不邪僻。知我是其道本所出，必不害我也。}

「庾公之斯至，曰：『夫子何爲不執弓？』曰：『今日我疾作，不可以執弓。』曰：『小人學射於尹公之他，尹公之他學射於夫子。我不忍以夫子之道反害夫子。雖然，今日之事，君事也，我不敢廢。』抽矢，叩輪，去其金，發乘矢而後反。」\jiazhu[1]{庾公之斯至竟，如孺子之所言。而曰我不敢廢君事，故叩輪去鏃，使不害人，乃以射孺子。禮射四發而去。乘，四也。《詩》云：「四矢反兮。」孟子言是以明羿之罪。假使如子濯孺子之得尹公之他而教之，何由有逢蒙之禍？}\par \jiazhu[1]{章指言：求交取友，必得其人。得善以全，養凶獲患。是故子濯濟難，夷羿以殘。可以鑒也。}

\section{論貌行善惡}

孟子曰：「西子蒙不潔，則人皆掩鼻而過之。」\jiazhu[1]{西子，古之好女西施也。蒙不潔，以不潔污巾帽而蒙其頭也。面雖好，以蒙不潔，人過之者皆掩鼻，懼聞其臭也。}

「雖有惡人，齋戒沐浴，則可以祀上帝。」\jiazhu[1]{惡人，醜類者也。面雖醜，而齋戒沐浴，自治潔淨，可以侍上帝之祀。言人當自治以仁義，乃爲善也。}\par \jiazhu[1]{章指言：貌好行惡，西子冒臭；醜人潔服，供事上帝。明當修飾，惟義爲常也。}

\section{論天下之言性}

孟子曰：「天下之言性也，則故而已矣。故者以利爲本。」\jiazhu[1]{言天下萬物之情性，當順其故則利之也。改戾其性，則失其利矣。若以杞柳爲桮棬，非杞柳之性也。}

「所惡於智者，爲其鑿也。如智者若禹之行水也，則無惡於智矣。禹之行水也，行其所無事也。」\jiazhu[1]{惡人欲用智而妄穿鑿，不順物。禹之用智，決江疏河，因水之性，因地之宜，引之就下，行其空虛無事之處。}

「如智者亦行其所無事，則智亦大矣。」\jiazhu[1]{如用智者不妄改作，作事循理，若禹行水於無事之處，則爲大智也。}

「天之高也，星辰之遠也，苟求其故，千歲之日至，可坐而致也。」\jiazhu[1]{天雖高，星辰雖遠，誠能推求其故常之行，千歲日至之日，可坐知也。星辰日月之會。致，至也，知其日至在何日也。}\par \jiazhu[1]{章指言：能修性守故，天道可知；妄智改常，必與道乖，性命之旨也。}

\section{論朝廷之禮}

公行子有子之喪，右師往弔。入門，有進而與右師言者，有就右師之位而與右師言者。\jiazhu[1]{公行子，齊大夫也。右師，齊貴臣王驩，字子敖。公行之喪，齊卿大夫以君命會，各有位次，故下云朝廷也。與言者，皆諂於貴人也。}

孟子不與右師言。右師不悅曰：「諸君子皆與驩言，孟子獨不與驩言，是簡驩也。」\jiazhu[1]{右師謂孟子簡其無德，故不與言，是以不悅也。}

孟子聞之，曰：「禮，朝廷不歷位而相與言，不踰階而相揖也。我欲行禮，子敖以我爲簡，不亦異乎？」\jiazhu[1]{孟子聞子敖之言，曰我欲行禮，故不歷位而言，反以我爲簡，異也。云以禮者，心惡子敖而外順其辭也。}\par \jiazhu[1]{章指言：循禮而動，不合時人。阿意事貴，脅肩所尊，俗之情也。是以萬物皆流，而金石獨止。}

\section{論君子存心}

孟子曰：「君子所以異於人者，以其存心也。君子以仁存心，以禮存心。仁者愛人，有禮者敬人。愛人者人恒愛之，敬人者人恒敬之。」\jiazhu[1]{存，在也。君子之在心者，仁與禮也。愛敬施行於人，人必反之己也。}

「有人於此，其待我以橫逆，則君子必自反也：我必不仁也，必無禮也，此物奚宜至哉？」\jiazhu[1]{橫逆者，以暴虐之道來加我也。君子反自思省，謂己仁禮不至也。物，事也。推此人何爲以此事來加我？}

「其自反而仁矣，自反而有禮矣，其橫逆由是也，君子必自反也：我必不忠。」\jiazhu[1]{君子自謂我必不忠。}

「自反而忠矣，其橫逆由是也，君子曰：『此亦妄人也已矣。如此，則與禽獸奚擇哉？於禽獸又何難焉？』」\jiazhu[1]{妄人，妄作之人，無知者。與禽獸何擇異也？無異於禽獸，又何足難也？}

「是故君子有終身之憂，無一朝之患也。乃若所憂則有之：舜，人也；我，亦人也。舜爲法於天下，可傳於後世，我由未免爲鄉人也，是則可憂也。」\jiazhu[1]{君子之憂，憂不如堯舜也。}

「憂之如何？如舜而已矣。」\jiazhu[1]{憂之當如之何乎？如舜而後可，故終身憂也。}

「若夫君子所患則亡矣。非仁無爲也，非禮無行也。如有一朝之患，則君子不患矣。」\jiazhu[1]{君子之行，本自不致意，常行仁行禮。如有一朝橫來之患，非己愆也，故君子歸天，不以爲患也。}\par \jiazhu[1]{章指言：君子責己，小人不改，比之禽獸，不足難矣。蹈仁行禮，不患其患，惟不若舜，可以憂也。}

\section{論禹稷顏回同道}

禹、稷當平世，三過其門而不入，孔子賢之。\jiazhu[1]{當平世三過其門者，身爲公卿，憂民急也。}

顏子當亂世，居於陋巷，一簞食，一瓢飲，人不堪其憂，顏子不改其樂，孔子賢之。\jiazhu[1]{當亂世安陋巷者，不用於世，窮而樂道也。}

孟子曰：「禹、稷、顏回同道。」\jiazhu[1]{孟子以爲憂民之道同，用與不用之宜若是也，故孔子俱賢之。}

「禹思天下有溺者，由己溺之也；稷思天下有飢者，由己飢之也，是以如是其急也。」\jiazhu[1]{由，猶也。禹、稷急民之難若是。}

「禹、稷、顏子易地則皆然。」\jiazhu[1]{顏子與之易地，其心亦然。不在其位，勞佚異矣。}

「今有同室之人鬪者，救之，雖被髮纓冠而救之，可也；鄉鄰有鬪者，被髮纓冠而往救之，則惑也，雖閉戶可也。」\jiazhu[1]{纓冠者，以冠纓貫頭也。鄉鄰，同鄉也。同室相救，是其理也，喻禹、稷走赴；鄉人非其事，顏子所以闔戶而高枕也。}\par \jiazhu[1]{章指言：上賢之士，得聖一槩。顏子之心，有同禹、稷。時行則行，時止則止，失其節則惑矣。}

\section{論匡章不孝}

公都子曰：「匡章，通國皆稱不孝焉。夫子與之遊，又從而禮貌之，敢問何也？」\jiazhu[1]{匡章，齊人也。一國皆稱不孝，問孟子何爲與之遊，又禮之以顏色喜悅之貌也？}

孟子曰：「世俗所謂不孝者五：惰其四支，不顧父母之養，一不孝也；博弈好飲酒，不顧父母之養，二不孝也；好貨財，私妻子，不顧父母之養，三不孝也；從耳目之欲，以爲父母戮，四不孝也；好勇鬪很，以危父母，五不孝也。章子有一於是乎？」\jiazhu[1]{惰，解不作。極耳目之欲，以陷罪戮及父母。凡此五者，人所謂不孝之行。章子豈有一事於是五不孝中也？}

「夫章子，子父責善而不相遇也。責善，朋友之道也；父子責善，賊恩之大者。」\jiazhu[1]{遇，得也。章子子父親教，相責以善，不能相得，父逐之也。朋友切磋，乃當責善耳。父子相責以善，賊恩之大也。}

「夫章子，豈不欲有夫妻子母之屬哉？爲得罪於父，不得近，出妻屏子，終身不養焉。」\jiazhu[1]{夫章子豈不欲身有夫妻之配，子有母子之屬哉？但以身得罪於父，不得近父，故出去其妻，屏遠其子，終身不爲妻子所養也。}

「其設心以爲不若是，是則罪之大者。是則章子已矣。」\jiazhu[1]{章子張設其心，執持此屏出妻子之意，以爲人得罪於父，而不若是以自責罰，是則罪益大矣。是章子之行已矣，何爲不可與言？}\par \jiazhu[1]{章指言：匡章得罪，出妻屏子，上不得養，下以責己。衆曰不孝，其實則否。是以孟子禮貌之也。}

\section{論曾子子思同道}

曾子居武城，有越寇。或曰：「寇至，盍去諸？」\jiazhu[1]{盍，何不也。曾子居武城，有越寇將來。人曰：「寇方至，何不去之？」}

曰：「無寓人於我室，毀傷其薪木。」寇退，則曰：「修我牆屋，我將反。」\jiazhu[1]{寓，寄也。曾子欲去，戒其守人曰：「無寄人於我室，恐其傷我薪草樹木也。」寇退則曰：「治牆屋之壤者，我將來反。」}

寇退，曾子反。左右曰：「待先生如此其忠且敬也。寇至則先去，以爲民望；寇退則反，殆於不可。」\jiazhu[1]{左右相與非議曾子者，言武城邑大夫敬曾子，武城人爲曾子忠謀，勸使避寇。君臣忠敬如此，而先生寇至則先去，使百姓瞻望而效之；寇退安寧則復來還，殆不可如是。怪曾子何以行之也。}

沈猶行曰：「是非汝所知也。昔沈猶有負芻之禍，從先生者七十人，未有與焉。」\jiazhu[1]{沈猶行，曾子弟子也。行謂左右之人曰：「先生之行非汝所能知也。先生，曾子也。往者先生嘗從門徒七十人舍吾沈猶氏，時有作亂者曰負芻，來攻沈猶氏，先生率弟子去之，不與其難。」言師賓不與臣同。}

「子思居於衛，有齊寇。或曰：『寇至，盍去諸？』子思曰：『如伋去，君誰與守？』」\jiazhu[1]{伋，子思名也。子思欲助衛君赴難。}

孟子曰：「曾子、子思同道。曾子，師也，父兄也；子思，臣也，微也。曾子、子思易地則皆然。」\jiazhu[1]{孟子以爲二人同道。曾子爲武城人作師，則其父兄，故去留無毀；子思微少也，又爲臣，委質爲臣當死難，故不去也。子思與曾子易處同然。}\par \jiazhu[1]{章指言：臣當營君，師有餘裕。二人處義非殊者，是故孟子紀之，謂得其同。}

\section{論堯舜與人同}

儲子曰：「王使人瞯夫子，果有以異於人乎？」\jiazhu[1]{儲子，齊人也。瞯，視也。果，能也。謂孟子曰：「王言賢者身貌必當有異，故使人視夫子，能有異於衆人之容乎？」}

孟子曰：「何以異於人哉？堯舜與人同耳。」\jiazhu[1]{人生同受法於天地之形，我當何以異於人哉？且堯舜之貌與凡同耳，其所以異，乃以仁義之道在於內也。}\par \jiazhu[1]{章指言：人以道殊，賢愚體別，頭圓足方，善惡如一。儲子之言，齊王之不達也。}

\section{齊人有一妻一妾}

齊人有一妻一妾而處室者。其良人出，則必饜酒肉而後反。其妻問所與飲食者，則盡富貴也。\jiazhu[1]{良人，夫也。盡富貴者，夫詐言其姓名也。}

其妻告其妾曰：「良人出，則必饜酒肉而後反；問其與飲食者，盡富貴也，而未嘗有顯者來。吾將瞯良人之所之也。」\jiazhu[1]{妻疑其詐，故欲視其所之。}

蚤起，施從良人之所之。遍國中無與立談者。卒之東郭墦間，之祭者乞其餘；不足，又顧而之他。此其爲饜足之道也。\jiazhu[1]{施者，邪施而行，不欲使良人覺也。墦間，郭外冢間也。乞其祭者所餘酒肉也。}

其妻歸，告其妾曰：「良人者，所仰望而終身也。今若此！」與其妾訕其良人，而相泣於中庭。\jiazhu[1]{妻妾於中庭悲傷其良人，相對涕泣而謗毀之。}

而良人未之知也，施施從外來，驕其妻妾。\jiazhu[1]{施施猶扁扁，喜悅之貌，以爲妻妾不知，如故驕之也。}

由君子觀之，則人之所以求富貴利達者，其妻妾不羞也而不相泣者，幾希矣。\jiazhu[1]{由，用也。用君子之道觀今求富貴者，皆以枉曲之道，昏夜乞哀而求之，以驕人於白日。由此良人爲妻妾所羞，爲所泣傷也。幾希者，言今苟求富貴，妻妾雖不羞泣者，與此良人妻妾何異也？}\par \jiazhu[1]{章指言：小人苟得，謂不見知；君子觀之，與正道乖。妻妾猶羞，況於國人？著以爲戒，恥之甚焉。}
